\chapter{Datenvorbereitung und Evaluation}
\label{chap:datenvorbereitung}

Bevor ein Modell trainiert werden kann, müssen die Daten vorbereitet werden. Ausserdem brauchen wir Werkzeuge, um die Qualität eines Modells zu messen. Die häufigsten Schritte in der Datenvorbereitung sind:

\begin{itemize}
    \item \textbf{Train-Test-Split}: Aufteilung der Daten in Trainings- und Testdaten.
    \item \textbf{Features und Label}: Definition der Eingabedaten (Features) und der Zielvariable (Label).
    \item \textbf{Normalisierung}: Skalierung der Features, damit sie vergleichbar sind.
    \item \textbf{Gütemasse}: Auswahl von Metriken, um die Leistung des Modells zu bewerten.
\end{itemize}

All diese Schritte werden in diesem Kapitel behandelt. 

\section{Trainings- und Testdaten}

Selbstlernende Modelle müssen auf Daten trainiert werden, um Muster zu erkennen. Um die Leistung eines Modells realistisch zu bewerten, ist es wichtig, dass es auf \textbf{ungesehenen} Daten getestet wird. Daher teilen wir die verfügbaren Daten in zwei Teile auf: den \textbf{Trainingsdatensatz} und den \textbf{Testdatensatz}.

\begin{mydefinition}[Train-Test-Split]
Die Daten werden in einen \textbf{Trainings\-datensatz} und einen
\textbf{Testdatensatz} aufgeteilt (typisch: 80\,\% / 20\,\%).
Das Modell wird \emph{nur} mit den Trainingsdaten trainiert und
anschliessend auf den \emph{ungesehenen} Testdaten bewertet.
\end{mydefinition}

\begin{mycaution}
Wer das Modell auch auf den Trainingsdaten bewertet, misst lediglich, wie gut das Modell \enquote{auswendig gelernt} hat -- das gibt keinen Aufschluss über die echte Leistung auf neuen Daten.
\end{mycaution}

\begin{myexercise}[Train-Test-Split nachvollziehen]
Ein Datensatz enthält 500 Zeilen.
\begin{enumerate}
    \item Wie viele Zeilen enthält der Trainings- bzw. der Testdatensatz bei einem 80/20-Split?
    \item Warum darf man die Testdaten \emph{nicht} für die Modellentwicklung verwenden?
    \item Was könnte in diesem Kontext mit \emph{data leakage} gemeint sein und warum wäre dies problematisch?
\end{enumerate}
\begin{myanswer}
\begin{enumerate}
    \item 400 Training / 100 Test \quad
    \item Sonst misst man nur, ob das Modell die Testdaten auswendig kennt, nicht ob es generalisiert. \quad
    \item Information aus dem Testset fliesst in das Training -- der gemessene Fehler unterschätzt den echten Fehler.
\end{enumerate}
\end{myanswer}
\end{myexercise}

\section{Merkmale und Zielvariable}
Im maschinellen Lernen werden die Eingabedaten als \textbf{Features} (Merkmale) bezeichnet, während die vorherzusagende Variable als \textbf{Label} (Zielvariable) bezeichnet wird.
\begin{mydefinition}[Features und Label]
    Essentielle Begriffe:
\begin{itemize}
    \item \textbf{Features} (Merkmale): die Eingabevariablen $x_1, x_2, \ldots, x_n$
    \item \textbf{Label} (Zielvariable): der Ausgabewert $y$, den das Modell vorhersagen soll
\end{itemize}
\end{mydefinition}

\begin{myexample}
Beim Datensatz \texttt{learning\_style\_classifier.csv}:
\begin{itemize}
    \item Features: Lernzeit, Schlafstunden, Anzahl Wiederholungen
    \item Label: Bestanden (1) oder nicht bestanden (0)
\end{itemize}
\end{myexample}

\section{Normalisierung}

\begin{mydefinition}[Min-Max-Normalisierung]
Features mit sehr unterschiedlichen Wertebereichen können Modelle verzerren.
Die \textbf{Min-Max-Normalisierung} skaliert alle Werte auf $[0,\,1]$:
\[
x' = \frac{x - x_{\min}}{x_{\max} - x_{\min}}
\]
Dies ist besonders wichtig für Algorithmen, die auf Abständen basieren (z.\,B. kNN, SVM). Ohne Normalisierung könnte ein Feature mit grösseren Werten die Distanzberechnung dominieren und somit die Modellleistung schmälern.
\end{mydefinition}

\begin{myexercise}[Normalisierung berechnen]
Gegeben: Lernzeiten $\{2, 4, 6, 8, 10\}$ Stunden.
\begin{enumerate}
    \item Normalisieren Sie alle Werte mit Min-Max. Sie können dazu das Python-Notebook verwenden: \href{\notebookbaseurl00_Datenvorbereitung.ipynb}{\texttt{00\_Datenvorbereitung.ipynb}}.
    \item Welcher Wert entspricht nach der Normalisierung dem Wert 6?
\end{enumerate}
\begin{myanswer}
$x_{\min}=2,\ x_{\max}=10$: 
$\{0,\; 0.25,\; 0.5,\; 0.75,\; 1.0\}$. 
Der Wert 6 wird zu $0.5$.
\end{myanswer}
\end{myexercise}

\section{Gütemasse für Klassifikation}

\begin{mydefinition}[Konfusionsmatrix]
Die \textbf{Konfusionsmatrix} zeigt, wie oft das Modell richtig und falsch klassifiziert hat:

\begin{center}
\begin{tabular}{cc|cc}
 & & \multicolumn{2}{c}{\textbf{Vorhergesagt}} \\
 & & Positiv & Negativ \\
\hline
\multirow{2}{*}{\textbf{Tatsächlich}} & Positiv & TP & FN \\
 & Negativ & FP & TN \\
\end{tabular}
\end{center}
\end{mydefinition}

\begin{mydefinition}[Accuracy, Precision, Recall]
\[
\text{Accuracy} = \frac{TP + TN}{TP + TN + FP + FN}
\qquad
\text{Precision} = \frac{TP}{TP + FP}
\qquad
\text{Recall} = \frac{TP}{TP + FN}
\]

Die \textbf{Accuracy} misst die Gesamtgenauigkeit des Modells.
Die \textbf{Precision} misst, wie viele der als positiv vorhergesagten Fälle tatsächlich positiv sind.
Der \textbf{Recall} misst, wie viele der tatsächlich positiven Fälle korrekt erkannt wurden.

Die Begriffe \emph{Sensitivität} und \emph{Spezifizität} werden manchmal synonym zu Recall bzw. Precision verwendet.
\end{mydefinition}

\begin{myexercise}[Metriken berechnen]
Ein Spam-Filter klassifiziert 200 E-Mails.
Die Konfusionsmatrix ergibt: 
\begin{center}
\begin{tabular}{cc|cc}
 & & \multicolumn{2}{c}{\textbf{Vorhergesagt}} \\
 & & Spam & Kein Spam \\
\hline
\multirow{2}{*}{\textbf{Tatsächlich}} & Spam & 80 & 20 \\
 & Kein Spam & 10 & 90 \\
\end{tabular}
\end{center}

\begin{enumerate}
    \item Berechnen Sie Accuracy, Precision und Recall. Sie können dazu das Python-Notebook verwenden: \href{\notebookbaseurl00_Datenvorbereitung.ipynb}{\texttt{00\_Datenvorbereitung.ipynb}}.
    \item Warum ist Recall besonders wichtig, wenn es darum geht, keinen Spam zu verpassen?
\end{enumerate}
\begin{myanswer}
\begin{enumerate}
    \item Berechnungen:
    \begin{itemize}
        \item Accuracy = $\frac{80 + 90}{200} = 0.85$
        \item Precision = $\frac{80}{80 + 10} = 0.89$
        \item Recall = $\frac{80}{80 + 20} = 0.80$
    \end{itemize}
    \item Ein hoher Recall bedeutet, dass die meisten Spam-E-Mails korrekt erkannt werden. Wenn der Recall niedrig ist, könnte der Filter viele Spam-Nachrichten durchlassen, was unerwünscht ist.
\end{enumerate}
\end{myanswer}
\end{myexercise}